\subsection{Элементы нелинейного анализа} % §13

Пусть $f: \R^1 \to R^1$. Тогда можно определить $\Delta f = A \Delta x + o(\Delta x)$, где $A = f'(x_0)$.
% не будем в этом параграфе отвлекаться на комплексности

% метрическое пространство придумали между титаником и пулеметом максим
\begin{definition}
    Пусть $E_1, E_2$~--- вещественные банаховы пространства. $D \subset E_1$~--- область, $F: D \to E_2$~--- отображение. $F$ называется дифференциируемым по Фреше в точке $x_0$, 
    если существует $A \in \L(E_1, E_2)$ такой, что $\Delta F = F(x_0 + h) - F(x_0) = A h + \varepsilon(h)$, где $\frac{\|\varepsilon(h) \|}{\|h\|} \to 0$, при $h \to 0$. $A$ называют производной Фреше.
\end{definition}

\begin{definition}
    Дифференциалом Фреше в точке $h$ называется 
    $dF(x_0, h) = F'(x_0) h$, $F' \in \L(E_1, E_2), h \in E_1$.
\end{definition}

\begin{lemmanote}
    Производная~--- это оператор, дифференциал~--- это значение этого оператора на элементе приращения.
\end{lemmanote}
% дуэль Коновалова и Редкозубова

\begin{definition}[эквивалентное определение]
    \[
    \frac{\|\Delta F - Ah\|}{\|h\|} \to 0
    \]
\end{definition}

\begin{statement}
    \begin{itemize}
        \item[0.]  Дифференцируемость по Фреше влечет непрерывность в точке.
        \item[1.] Если $F(x) = \text{const} (\in E_2) \forall x \in D$, то $\forall x \in D F'(x) = 0$.
        \item[2.] $F \in \L(E_1,E_2)$. $\forall x \in E_1 \; F'(x)=F$
        \item[3.] $(\alpha_1 F_1 + \alpha_2 F_2)'=\alpha_1 F_1' + \alpha_2 F_2'$
        \item[4.] $F: E_1 \to E_2, G: E_2 \to E_3$~--- дифференцируемы. 
        Тогда $H = G \circ F$ дифференцируема и $H' = G' \circ F'$. % Упражнение: доказать 4 пункт
    \end{itemize}
\end{statement}

% ГОСТ 7.0.12-2011 Правильно сокращать дифференцируемость как диф., а не дифф. Но мы будем не по ГОСТу

\begin{definition}
    $F: D \to E_2$. Если существует $DF(x_0; h) = \left/ \frac{d}{dt} F(x_0 + th) \right|_{t = 0}$, то говорим, что определен дифференциал Гато.
    $DF(x_0; h) = Ah$, где $A = F'_{\text{сл.}}(x_0)$ называется слабой производной Гато.
\end{definition}
\begin{statement}
    \begin{itemize}
        \item[1.] Дифференцируемость по Фреше $\implies$ дифференцируемость по Гато и $F'_{\text{сл.}}=F'_{\text{Фр.}}$.
        \item[2.] В обратную сторону это неверно
        \textit{Картинка: на всей плоскости ноль, на параболе (кроме точки ноль) -- 1}
    \end{itemize}
\end{statement}

\begin{example}
В пространстве $C[a,b]$
\[
(Af)(x) = \int_a^b K(x, t, f(t)) dt
\]

\[
\frac{d}{dx} \int_a^b K(x, t, f(t)) dt
\]
Когда дифференцирование и интеграл можно переставить местами?
\end{example}

\begin{exercise}
    $H(\R), f(x)=\|x\|$. Где он дифференцируем? (ответ: везде кроме 0) Найти производную.
\end{exercise}

% \begin{theorem}[13.1 (Лагранжа о конечных приращениях)]\label{13.1}
%     Пусть $E_1, E_2$~--- вещественные банаховы пространства,
%     $D \subset E_1$~--- выпуклая область, $F: D \to E_2$~--- дифференцируема (по Фреше). Тогда $\forall x_0, x_1 \in D$ выполнено $\|F(x_1) - F(x_0)\| \leqslant \sup_{y \in (x_0, x_1)} \|F'(y)\| \cdot \|x_1 - x_0\|$. 
% \end{theorem}

% \begin{proof}

% % TODO картинка сэндвич
% % здесь история, которую вырежут на записи

% $f \circ F \circ x(t)$ дифференцируема.

% \end{proof}